\section{Two orbits of decorated tetrahedra}

Write $x=[(1,0)^T]_K$, $y=[(0,1)^T]_K$ and $z=[(1,1)^T]_K$.
Put
\[
 T=\{x,y,z\},\qquad d_2=[(-e^2,1)^T]_K,\quad
 d_3=[(-e,1)^T]_K,\qquad D_i=T\cup\{d_i\}.
\]
The four points of either $D_i$ have different projective lines.
Every determinant of two distinct representatives in $D_i$ belongs
to $K$. For example the determinants involving $(1,1)^T$ and
$(-e^2,1)^T$ give $1+e^2=-e$; the other five determinants are
$\pm1$ or $\pm e^2$. The calculation for $D_3$ exchanges $e,e^2$.
Let $\mathcal D_i$ be the $G$-orbit of $D_i$.

\begin{lemma}\label{tetra-stabilizers}
The stabilizers of $D_2,D_3$ are respectively
$G_2=\gen{a_1,a_3}$ and $G_3=\gen{a_1,a_2}$, each acting as
$A_4$ on its four decorated points. Suppose $p$ is odd.
Then the stabilizer of $T$ is $\gen{a_1}\cong C_3$,
the two tetrahedron orbits are disjoint, and each triple in a
tetrahedron of one orbit belongs to exactly one tetrahedron of
the other orbit.
\end{lemma}
\begin{proof}
On either tetrahedron $a_1$ cycles $x,y,z$ and fixes its fourth
point. On $D_2$, $a_3$ fixes $x$ and cycles $y,z,d_2$;
on $D_3$, $a_2$ fixes $x$ and cycles $y,z,d_3$.
These assertions follow by multiplying the four column vectors;
all scalar factors are in $K$. The two displayed three-cycles
generate $A_4$. A projective linear transformation fixing three
different projective lines is the identity, so these actions are
faithful and both subgroups have order twelve.

Consider an element stabilizing $T$. If it fixes each of $x,y,z$,
a lift to $\operatorname{SL}_2(F)$ is diagonal, and fixing $z$
forces its diagonal entries to agree. It is scalar and hence the
identity in $G$. Thus the stabilizer embeds in $S_3$ and contains
the three-cycle $a_1$. An element transposing $x,y$ and fixing $z$
would have a lift
\[
                         \begin{pmatrix}0&k\\k&0\end{pmatrix},
                         \qquad k\in K,
\]
with determinant $-k^2=1$. For odd $p$, the square subgroup of
$K$ has order three and cannot contain $-1$. There is no such
transposition; composing any odd permutation with a power of
$a_1$ reduces to this case. The stabilizer is therefore $C_3$.

For $p$ odd, an element stabilizing $D_i$ can first be composed
with an element of the displayed $A_4$ to fix $d_i$.
It then stabilizes $T$ and already belongs to that $A_4$.
This proves the full stabilizer assertion. If a translate $gD_i$
contains $T$, transitivity of $G_i$ on the four triples of $D_i$
allows us to choose $h\in G_i$ with $ghT=T$. Consequently
$gD_i=ghD_i=D_i$. Thus $D_i$ is the unique member of its orbit
containing $T$. In particular $D_2,D_3$ are in different orbits.
Transporting this argument by $G$ proves the assertion for every
triple in either orbit.

The remaining characteristic is $p=2$, so $q=4$ and $K=F^\times$.
The decorated points are the five ordinary projective points.
The group $G$ has order sixty and acts faithfully on them, with
the three generators acting by even permutations. It is $A_5$.
Both $D_i$ are four-subsets, and their orbits are the five
four-subsets. Their stabilizers have order twelve and are precisely
the already exhibited $G_i$. The assertions specific to odd $p$
are not used in this case.
\end{proof}

Define a four-type incidence system as follows. Types $0,1$ are
two copies of the decorated points; types $2,3$ are the orbits
$\mathcal D_2,\mathcal D_3$, retaining the type even if the orbits
coincide. Between different types use
\[
\begin{array}{c|l}
\text{types}&\text{incidence}\\\hline
0,1&\det(v,w)\in K,\\
i,j\ (i<2\le j)&[v]_K\in D,\\
2,3&|D\cap D'|=3.
\end{array}                                                     \tag{5}
\]

\begin{theorem}\label{full-geometry}
The incidence system (5) is the source's right-coset system
$\Gamma(G;(G_i)_{i=0}^3)$. It is a hypertope with tetrahedral
diagram having all six bonds labelled three. The group $G$ has
two free orbits on its chambers; adjacent chambers belong to
different orbits. All rank-three flags form one free $G$-orbit
for each choice of three types.
\end{theorem}
\begin{proof}
The point stabilizers and the first incidence rule were proved in
Theorem~\ref{incidence}. Lemma~\ref{tetra-stabilizers} identifies
the other two stabilizers. Thus the object map sends a right coset
$G_i g$ to $g^{-1}$ applied to the corresponding base object
$y,x,D_2,D_3$. Every incident point--tetrahedron pair is in the
orbit of the corresponding base pair, because $A_4$ is transitive
on the four vertices. For odd $p$, every tetrahedron of type $2$
has exactly four neighbours of type $3$, one for each triple.
Its $A_4$ stabilizer is transitive on these triples. Hence these
pairs too form exactly the orbit of the base pair $(D_2,D_3)$.
Pair incidence in a coset system is exactly the orbit of its base
pair, proving the claim about all six relations.

For the moment suppose $p$ is odd. A flag of types $012$ or $013$
is a tetrahedron with two ordered distinct vertices. Its stabilizer
$A_4$ acts regularly on the twelve ordered pairs of distinct
vertices. These rank-three flags therefore form one free orbit.
A flag of types $023$ or $123$ consists of adjacent tetrahedra
and a point of their common triple. The stabilizer of the pair is
$C_3$, acting regularly on that triple. These flags likewise form
one free orbit.

There are exactly two extensions of each such flag to a chamber.
For types $023$ or $123$, the missing point can be either of the
other two vertices in the common triple. For types $012$ or
$013$, the missing tetrahedron is specified by a triple of the
given tetrahedron containing the two given points. There are
exactly two choices, giving different opposite tetrahedra.
Every incident pair extends to one of these flags: for a point
and a tetrahedron choose another vertex; for adjacent tetrahedra
choose a point in their common triple; and the point pair case
follows by transporting the base point pair. Every object has
an incident pair. Thus this is a geometry, and the two-extension
calculation proves thinness.

The stabilizer of each base incident pair is cyclic of order
three. For pair types $01$, $02$, $03$, $12$, $13$ and $23$,
these groups are respectively
\[
 \gen E,\quad\gen{a_1^{-1}a_3},\quad\gen{a_1^{-1}a_2},\quad
 \gen{a_3},\quad\gen{a_2},\quad\gen{a_1}.
\]
For point--tetrahedron pairs this follows from the $A_4$ action;
the other two cases were proved above. Since triple stabilizers
are trivial, every rank-two residue has three objects of each
of its two types. Thinness makes its graph a simple bipartite
graph of degree two, hence a six-cycle. This proves both
connectedness of these residues and the bond labels.

Each rank-three residue is the coset geometry inside its object
stabilizer $G_i$ with parabolics $G_i\cap G_j$. This assertion
includes incidence: a pair in such a residue, together with the
fixed object, is a rank-three flag, so the transitivity just
proved puts it in the orbit of the base flag under $G_i$.
These three cyclic parabolics generate $G_i$ in every case, as
the displayed generators show. Their coset incidence graph is
connected: writing a group element as a product of elements of
these parabolics gives a path between the corresponding chambers
and their objects. The whole geometry is connected already on
types $0,1$, and every tetrahedron has point neighbours. All
required residues are therefore connected.

There are exactly $2|G|$ chambers, because each rank-three flag
has two extensions and there are $|G|$ flags of any fixed three
types. Chamber stabilizers are trivial. If adjacent chambers
were in the same orbit, a group element taking one to the other
would fix their common rank-three flag and would be the identity,
a contradiction. This proves the orbit assertions.

For $q=4$, identify each tetrahedron with its missing projective
point. All four types are then copies of the same five-element
set, with incidence meaning unequal labels. A chamber is an
ordered choice of four different labels. Every three-label flag
has two extensions; every rank-two residue is a six-cycle, and
every rank-three residue is three copies of four labels with
unequal-label incidence, which is connected. The $A_5$ action
on ordered triples of different labels is free and transitive.
These observations give all the same conclusions directly,
including the coset identification by the stabilizers.
\end{proof}

The six pair stabilizers are distinct, and their intersections
are trivial whenever a third base object is fixed. Together with
$G_i\cap G_j$ as displayed, this proves the full intersection
property for the rotation marking: a subgroup indexed by two
labels is one of these cyclic groups, one indexed by three is a
maximal parabolic, and the intersections are their indexed meets.
For example intersecting a cyclic group with a maximal parabolic
that does not contain it fixes three base objects and is trivial.
Two distinct cyclic groups have trivial intersection by the same
argument. The generators are independent since omitting one
leaves a proper maximal parabolic. This also recovers the source's
$C^+$-group claims, rather than assuming them to construct the
geometry.

\section{Residues and symmetries}

The residue of a tetrahedron is three copies of its four vertices:
the opposite tetrahedron is labelled by the one missing vertex
of its common triple. Incidence again means unequal labels.
Its rotation group is $A_4$; the two triangular translation words
are distinct nonidentity double transpositions and generate its
Klein four subgroup. Their squares are the identity, so the
translation kernel contains $2\Z^2$. Its index is four because
the image is that Klein four group. The kernel is therefore
exactly $2\Z^2$, giving the marked hypermap $(3,3,3)_{(0,2)}$.

The residue of a point has rotation group $G_0$ or $G_1$, with
the ordered generators used in the next section. Its actual
incidence is the internal coset incidence by the rank-three flag
argument in Theorem~\ref{full-geometry}. The classical triangular
rotation group maps onto it with a translation kernel. Therefore
the residue is the triangular toroidal quotient by precisely that
kernel, not merely by a relation satisfied in the finite group.
In either point residue the three cyclic parabolics, in order,
are $\gen{X^{-1}Y}$, $\gen Y$ and $\gen X$; they have order
three and trivial pairwise intersections. Thus this is the
standard marked $C^+$ coset construction. The residue stabilizer
has two free chamber orbits, inherited from the two global orbits,
and adjacent chambers are in opposite orbits.
The full kernels and the resulting correction of the printed
period parameters are calculated below.

Consider the two linear maps
\[
       A=\begin{pmatrix}-1&1\\0&1\end{pmatrix},\qquad
       B=\begin{pmatrix}0&1\\1&0\end{pmatrix}.
\]
Both have determinant $-1\in K$ and normalize $\operatorname{SL}_2(F)$.
They act on decorated points and tetrahedra, preserve all the
relations (5), and exchange $D_2,D_3$. The latter assertion
follows by applying each map to their four displayed vectors.
Thus $A$, accompanied by the type permutation $(2\ 3)$, is a
correlation $\mu$. It sends the base chamber
$(y,x,D_2,D_3)$ to $(z,x,D_2,D_3)$, its type-zero neighbour.
The map $B$, accompanied by $(0\ 1)(2\ 3)$, is a correlation
$\eta$ fixing the base chamber as a set. Both normalize the
$G$ action. Accordingly $\mu$ exchanges its two chamber orbits,
whereas $\eta$ preserves them. This proves the improper and
proper correlation assertions by the same object model.

When $q=p^2$, Frobenius $\sigma:t\mapsto t^p$ exchanges $e,e^2$.
On the object model it preserves determinant membership in $K$,
fixes $x,y,z$ and exchanges $D_2,D_3$. Accompanied by the type
permutation $(2\ 3)$ it is a correlation fixing the base chamber.
Its composition with $\mu$ is type preserving and takes the base
chamber to its type-zero neighbour. Together with $G$ it acts
transitively on all chambers. The hypertope is regular.

For $q=p$, the kernel calculation below shows that neither
marked point residue admits an automorphism inverting its two
standard rotation generators. The acquired theorem on two chamber
orbits~\cite{reflection}, applied with the preceding actual
coset and chamber identifications, says that the residue is
regular exactly when this inversion descends. Its proof labels
all chambers $(g,+),(g,-)$ and uses the explicit reflection
$(g,\pm)\mapsto(\theta(g),\mp)$ for an inverting automorphism
$\theta$; the converse follows by conjugating the rotations
by a normalized chamber reflection. The non-invariant kernel
therefore makes the point residue chiral.
A regular geometry has regular residues:
chamber transitivity restricts to chamber transitivity in the
stabilizer of the fixed flag. Hence the full geometry is not
regular. Its two $G$-orbits and thin residual connectedness show
that it is chiral, with full type-preserving group $G$.
Here we use the elementary gallery fact that in a thin residually
connected geometry any type-preserving automorphism fixing a
chamber is the identity: propagate across its uniquely determined
adjacent chambers. Residual connectedness gives gallery
connectedness by induction on rank. In particular the full group
cannot split either transitive $G$-orbit or contain an additional
element preserving both orbits.

For $q=4$ the full correlation group is $S_5\times S_4$.
Indeed the common permutation of the five labels and any
permutation of the four types preserve unequal-label incidence.
Conversely, remove the induced type permutation from a correlation.
For two different types, the unique non-neighbour of a label is
its equal label in the other copy. The permutations on the four
copies must therefore agree. The type-preserving group is $S_5$,
and the given rotation group is $A_5=\operatorname{PSL}_2(4)$.
This corrects the final isomorphism in the source's first remark
on page~19: $S_5\times S_4$ is not
$\operatorname{PSL}_2(4)\times S_4$.
